\section{Infinitely many failures in one fixed map}\label{sec:infinite}

The realization theorem changes the Keller map when the polynomial changes.
A stronger phenomenon occurs for one quadratic-gauge map.

Put
\[
 G_\infty(S)=S^5-\frac32S^4+\frac32S^3
             -\frac54S^2+\frac9{16}S
\]
and let $F_\infty=F_{G_\infty}$ be the map of
Proposition~\ref{prop:gauge}.  Thus $\det DF_\infty=-2$.

\begin{theorem}[One map, infinitely many Hasse failures]
For every rational prime $\ell\equiv1\pmod{27}$, put
\[
 y_\ell=
 \left(1,\frac{32\ell}{9},\frac{8\ell+1}{3}\right).
\]
Then
\[
 F_\infty^{-1}(y_\ell)
 \simeq
 \Spec\Q[X]/\bigl((X^3-\ell)(X^2+X+1)\bigr).
\]
Each fiber has points over every completion of $\Q$ but no rational point.
The targets are pairwise distinct, and with the standard projective height,
\[
 \#\{y_\ell:H(y_\ell)\leq B\}
 \sim\frac{B}{576\log B}.
\]
\end{theorem}

\begin{proof}
Set $X=S-\tfrac12$.  The identity
\[
 (S-\tfrac12)^3(S^2+\tfrac34)=G_\infty(S)-\frac3{32}
\]
shows that the inverse equation at $y_\ell$ is exactly
\[
 (X^3-\ell)(X^2+X+1).
\]
It is squarefree and has no rational root.  The local argument is the same as
in Theorem~\ref{thm:hasse}.  At $p=3$, the stronger congruence
$\ell\equiv1\pmod{27}$ gives a cubic root by Hensel's lemma; at $p=\ell$, the
quadratic factor splits because $\ell\equiv1\pmod3$.  All other primes are
covered by the same residue-class dichotomy.  Dirichlet's theorem gives
infinitely many such primes.

The primitive projective coordinates of $y_\ell$ are
$[9:9:32\ell:24\ell+3]$, so $H(y_\ell)=32\ell$.  The prime number theorem in
the progression $1\bmod27$ gives the displayed asymptotic.
\end{proof}
